\section{Einleitung}

Die vorliegende Projektarbeit beschreibt die Entwicklung einer analytischen Datenanwendung zur Verarbeitung und Bewertung von Insider‑Transaktionsdaten. Öffentliche Meldungen über Aktienkäufe und ‑verkäufe durch Führungskräfte gelten als Indikatoren für die Erwartungen dieser Personen an den künftigen Unternehmenserfolg. Die United States Securities and Exchange Commission (SEC) verpflichtet Insider, ihre Transaktionen über das Formular~4 offenzulegen. Diese Meldungen sind über verschiedene Programmierschnittstellen verfügbar.

Mercator nutzt die Schnittstelle von Financial Modeling Prep (FMP), um aktuelle Insider‑Transaktionen abzurufen und mit Markt‑ und Unternehmensdaten anzureichern. Im Rahmen des Moduls \emph{Datenbanken 2} wird eine Anwendung entwickelt, die einen vollständigen Datenfluss von der Erfassung über die Bereinigung bis hin zur Speicherung und visuellen Analyse realisiert. Neben der technischen Umsetzung ist eine wissenschaftliche Reflexion erforderlich, die die theoretischen Grundlagen, die Anforderungen, die Architektur, die Implementierung und die Bewertung der Ergebnisse umfasst.

Die Einleitung stellt den Kontext der Arbeit dar, formuliert das Ziel und definiert die Forschungsfrage. Anschließend wird die Gliederung erläutert.

\subsection{Zielsetzung}

Ziel dieser Arbeit ist es, eine modulkonforme Lösung zu entwickeln, die öffentliche Insider‑Transaktionsdaten importiert, validiert, filtert, anreichert, bewertet, in zwei unterschiedlichen Datenbanksystemen speichert und über eine interaktive Benutzeroberfläche zugänglich macht. Die Anwendung soll sowohl die technische Machbarkeit als auch die analytische Nutzbarkeit demonstrieren. Die Ergebnisse dienen ausschließlich der technischen Evaluation und stellen keine Anlageempfehlung dar.

\subsection{Forschungsfrage}

Die Forschungsfrage lautet:

\begin{quote}
  Wie lässt sich eine analytische Datenanwendung zur Verarbeitung, Speicherung, Bewertung und Visualisierung öffentlicher Insider‑Transaktionsdaten mit relationaler und dokumentenorientierter Datenhaltung umsetzen?
\end{quote}

\subsection{Aufbau der Arbeit}

Kapitel~\ref{sec:grundlagen} behandelt die theoretischen Grundlagen von Datenverarbeitung, Datenbanken und Insider‑Transaktionsdaten. Kapitel~\ref{sec:anforderungen} beschreibt die funktionalen und nicht‑funktionalen Anforderungen. Kapitel~\ref{sec:architektur} stellt die Architektur und das Datenmodell vor. Kapitel~\ref{sec:umsetzung} erläutert die Implementierung der Datenpipeline und des Scorings. Kapitel~\ref{sec:anwendung} zeigt die Benutzeroberfläche. Kapitel~\ref{sec:bewertung} bewertet die Ergebnisse und diskutiert Grenzen. Kapitel~\ref{sec:fazit} schließt die Arbeit mit einem Fazit ab.